\scanpage{8}\mspage{7}
Je voudrais maintenant faire passer les axiomes de disjonction-compatibilité avant ceux sur les lieux, pour ne pas \tikz[baseline=(g.base),every node/.style={inner sep=0pt}]{\node (g) at (0,0) {}; \node[anchor=base west] (a) at (.1em,1.8ex) {\inserted{avoir l'air de}}; \draw (0,-.15ex) .. controls (.2em,1.1ex) and ([xshift=-.6em,yshift=-.5ex]a.base east) .. ([xshift=.1em,yshift=-.35ex]a.base east);} les en faire dépendre. Les axiomes sont ``facultatifs'' pour un atelier, et il y aura pas mal \unsure{des'|des|d'} implications mutuelles avec les axiomes sur les lieux.

Une notion essentielle ici est celle de ``$F$ et $G$ \uline{disjoints}''
\[
 F\parallel G\ \overset{\mathrm{dfn}}{\Longleftrightarrow}\ F\mathrel{\tikz[baseline=-.4ex,x=.65em,y=.65em]{\draw (0,1.1)--(1,.5)--(0,-.1) (1,1.1)--(0,.5)--(1,-.1);}}G\quad\text{et}\quad F\cap G=\varnothing_F
\]
notion qui ne dépend que de $\leq$. On a d'ailleurs
\[
 F\parallel G\ \Longleftrightarrow\ \forall X\in\widetilde F,\ Y\in\widetilde G,\quad X\parallel Y
\]
donc interprétation en \unsure{termes|terme} de $X\parallel Y$ dans $M$, qui y signifie
\[
\begin{aligned}
X\mathrel{\underset{M}{\parallel}}Y\ \Longleftrightarrow\
 X\mathrel{\underset{M}{\mathrel{\tikz[baseline=-.4ex,x=.65em,y=.65em]{\draw (0,1.1)--(1,.5)--(0,-.1) (1,1.1)--(0,.5)--(1,-.1);}}}}Y\quad &\text{et}\quad\{X,Y\}\ \text{n'est \uline{pas}}\\
&\text{minoré dans } M,\leq\\
&\text{i.e. }\nexists Z\ \text{avec }Z\leq X,Z\leq Y
\end{aligned}
\]

% The shared stem is not duplicated. The overwrite touches the t junction and ending;
% the suffix segmentation below is an approximation, with both layers uncertain.
(\uline{NB} Il résult\deleted{\unsure{era}}\inserted{\unsure{e|era}} des axiomes \tikz[baseline=(g.base),every node/.style={inner sep=0pt}]{\node (g) at (0,0) {}; \node[anchor=base west] (a) at (.1em,1.8ex) {\inserted{précédents}}; \draw (0,-.15ex) .. controls (.2em,1.1ex) and ([xshift=-.6em,yshift=-.5ex]a.base east) .. ([xshift=.1em,yshift=-.35ex]a.base east);} que ça implique que $\nexists Z$ avec $Z\ll X,Z\ll Y$\unsure{)|,} \deleted{i.e\unsure{.} c'est un $\unsure{\exists}$} et \deleted{\unsure{la|le}} réciproque résultera \unsure{d'axiome|d'axiomes} ci-dessous\ldots)

% The following At DC1/Cor draft is crossed by three long diagonal strokes.
% One of those strokes begins in, and crosses, the slanted left-margin note.
% Global cancellations are represented by deleted; nested deleted records
% the smaller local cancellations that predate the global abandonment.
\reversemarginpar
\margin{\deleted{Un atelier sera dit \uline{\unsure{spacieux|specieux}} s'il satisfait cette condition. Contre-ex. : les ateliers simpliciaux\ldots}}%
\normalmarginpar
\figure{\begin{tikzpicture}
\node[anchor=north west,inner sep=0pt,text width=.94\linewidth,font=\small] (dc) at (0,0) {%
\deleted{\uline{At DC} 1 \quad (\uline{Axiome des ateliers \unsure{spacieux|specieux}}) Soient $X,Y\in M$, \inserted{soit} \unsure{\tikz[baseline=-.25ex,x=.5em,y=.35em]{\draw (.13,.55)--(0,.2)--(1.35,.2)..controls(1.5,.22)and(1.29,.52)..(.98,.4)--(0,.2);}|\deleted{et}} $K=X\cap Y=\underset{Z\in\widetilde X\cap\widetilde Y}{\mathrm{Sup}}\,Z$. \deleted{\unsure{Il}} Supposons que \inserted{pour $X',Y'\in M$}}
\[
\deleted{\deleted{\forall}\ X'\ll X,\ X'\parallel K,\ Y'\ll Y,\ Y'\parallel K\ \Longrightarrow\ X'\parallel Y'.}
\]
\deleted{Alors [NB et alors seulement\ldots\
\mbox{\tikz[baseline=(b.base),every node/.style={inner sep=0pt}]{\node (b) at (0,0) {}; \node[anchor=base west] at (0,-1.6ex) {\inserted{cf Cor}}; \path (0,0)--(3.6em,0);}}
] $X\mathrel{\tikz[baseline=-.4ex,x=.65em,y=.65em]{\draw (0,1.1)--(1,.5)--(0,-.1) (1,1.1)--(0,.5)--(1,-.1);}}Y$.}

% The first overwritten index shows F on X, without a securely detached tilde.
% A single visible wave is above the second G/Y correction; its attribution
% to the old Y layer is unresolved, hence the old-layer alternative below.
\deleted{\uline{Cor} Soient $F,G\in\mathscr F$, $K=F\cap G=\underset{Z\in\deleted{\unsure{X|\widetilde X}}\inserted{F}\cap\deleted{\unsure{\widetilde Y|Y}}\inserted{\widetilde G}}{\mathrm{Sup}}\,Z$.}
\[
\begin{aligned}
&\deleted{\text{Alors } F\mathrel{\tikz[baseline=-.4ex,x=.65em,y=.65em]{\draw (0,1.1)--(1,.5)--(0,-.1) (1,1.1)--(0,.5)--(1,-.1);}}G\quad \text{\tikz[baseline=(s.base)]{\node[inner sep=1pt] (s) {ses};\draw (.55em,.62em)..controls(-.1em,.9em)and(-1.05em,.55em)..(-1.06em,-.08em)..controls(-1.07em,-.7em)and(.55em,-.92em)..(.99em,-.25em)..controls(1.2em,.17em)and(.94em,.46em)..(.75em,.5em);}}\quad \forall X,\deleted{\unsure{\varphi|y}}\inserted{Y}\in M}\\
&\deleted{X\ll F,\ Y\ll G,\ X\parallel K,\ Y\parallel K\ \Longrightarrow\ X\parallel Y}
\end{aligned}
\]
};
\draw ([xshift=.03\linewidth,yshift=1.4em]dc.north west)--([xshift=.31\linewidth,yshift=1em]dc.south west);
\draw ([xshift=.36\linewidth]dc.north west)--([xshift=.67\linewidth]dc.south west);
\draw ([xshift=.73\linewidth]dc.north west)..controls([xshift=-.03\linewidth]dc.east)and([xshift=.92\linewidth,yshift=4em]dc.south west)..([xshift=.94\linewidth,yshift=3em]dc.south west);
\end{tikzpicture}}
